\section{Réduction de Feshbach--Schur}

Soit

\[
\mathcal H=\mathcal H_0\oplus\mathcal H_{\rm ext},
\qquad
H=
\begin{pmatrix}
H_{00}&V\\ V^*&H_{\rm ext}
\end{pmatrix}.
\]

Pour \(z\) dans l'ensemble résolvant de \(H_{\rm ext}\), la compression de la
résolvante est

\begin{equation}
P_0(z-H)^{-1}P_0
=\left[z-H_{00}-\Sigma_0(z)\right]^{-1},
\qquad
\Sigma_0(z)=V(z-H_{\rm ext})^{-1}V^*.
\label{rm:eq:rm-feshbach-schur}
\end{equation}

L'opérateur \(\Sigma_0\) transporte vers le feuillet présent la mémoire des feuillets
elliptique et hyperbolique. Sa transformée temporelle engendre un noyau de
mémoire

\[
S_{\rm eff}[\psi_0]
=\int\psi_0^\dagger(i\hbar\partial_t-H_{00})\psi_0\,\mathrm dt
-\iint\psi_0^\dagger(t)K_0(t,t')\psi_0(t')\,\mathrm dt\,\mathrm dt'.
\]

Un développement régulier à basse fréquence,

\[
\Sigma_0(\omega)
=\Sigma_0^{(0)}+\omega\Sigma_0^{(1)}
+\omega^2\Sigma_0^{(2)}+\cdots,
\]

produit une action locale effective ; une structure spectrale singulière conserve
une mémoire non locale. Le présent reçoit ainsi la réponse de l'état global sans
introduire un temps extérieur.
